\section*{5 On the NP-hardness of communication minimization problems}
\subsection*{5.1 Background}
\paragraph{Hardness of graph coloring}
Our NP-hardness reduction is from the chromatic number problem:
\setcounter{nttheorem}{18}
\begin{ntdefinition}[Author Definition19: Chromatic number]
A coloring of an undirected graph $G$, is a partition of the vertices such that no edge has both endpoints in the same part. The chromatic number of a graph $G$, denoted $\chi(G)$, is the smallest number of parts of a coloring of $G$.
\end{ntdefinition}
The NP-hardness of approximating the chromatic number has been established by a series of results~\cite{ref57,ref31,ref24}, culminating in a paper by Zuckerman~\cite{ref82}, where the following was proven:
\begin{nttheorem}[Author Theorem20, {\cite{ref82}}]
For every constant $\epsilon>0$ it is NP-hard to approximate $\chi(G)$ for a given $n$-vertex graph $G$, with an approximation ratio better than $n^{1-\epsilon}$. More precisely, for every $L\in\operatorname{NP}$ and $\epsilon>0$, there exists a polynomial-time algorithm that, on input $x$, outputs a parameter $k$ and an $n$-vertex graph $G$ such that if $x\in L$ then $\chi(G)\leq k$, and if $x\notin L$ then $\chi(G)>n^{1-\epsilon}\cdot k$.
\end{nttheorem}
In the reductions above, the parameter $k$ is $\Theta(n^\epsilon)$. More recent results on the hardness of approximating chromatic number allow for a different gap, where $k$ is constant as $n$ grows, and we wish to distinguish graphs which are $k$-colorable from graphs which are not $g(k)$-colorable, for a fast-growing function $g:\mathbb N\to\mathbb N$. An original such result with $g(k)=k^{\Omega(\log k)}$ was shown by Khot~\cite{ref47}, which was later improved to $g(k)=2^{\Omega(k^{1/3})}$ by Huang~\cite{ref37}. The latest result, by Wrochna and \v Zivn\'y~\cite{ref79}, achieves $g(k)=\binom k{\lfloor k/2\rfloor}$:
\begin{nttheorem}[Author Theorem21, {\cite{ref79}}]
Let $k\geq4$ be a natural number. For every $L\in\operatorname{NP}$, there exists a polynomial-time algorithm that, on input $x$, outputs an graph $G$ such that if $x\in L$ then $\chi(G)\leq k$, and if $x\notin L$ then $\chi(G)>\binom k{\lfloor k/2\rfloor}$.
\end{nttheorem}
\paragraph{Communication complexity of relations and partial functions}
We will need the following definitions.
\begin{ntdefinition}[Author Definition22]
We will call a two-player communication relation, or simply a relation, to any subset $F\subseteq\mathcal X\times\mathcal Y\times\mathcal Z$, where $\mathcal X,\mathcal Y,\mathcal Z$ are finite sets, such that, for every $(x,y)\in\mathcal X\times\mathcal Y$, there exists at least one $z\in\mathcal Z$ with $(x,y,z)\in F$.
\end{ntdefinition}
Given a relation $F\subseteq\mathcal X\times\mathcal Y\times\mathcal Z$, and a pair $(x,y)\in\mathcal X\times\mathcal Y$, we let
\[F(x,y)=\{z\in\mathcal Z\mid(x,y,z)\in F\}\ne\varnothing.\]
The matrix with $\mathcal X$-indexed rows and $\mathcal Y$-indexed columns, whose $(x,y)$ entry is $F(x,y)$, is called the communication matrix of $F$.
Given a relation $F\subseteq\mathcal X\times\mathcal Y\times\mathcal Z$, a rectangle $R=A\times B\subseteq\mathcal X\times\mathcal Y$, and an element $z\in\mathcal Z$, we say $R$ is $z$-monochromatic for $F$, if $(x,y,z)\in F$ for all $(x,y)\in R$. We say $R$ is monochromatic for $F$ if there exists a $z\in\mathcal Z$ with $R$ being $z$-monochromatic for $F$.
A partial two-player function is a relation $f\subseteq\mathcal X\times\mathcal Y\times\mathcal Z$ such that, for every $(x,y)\in\mathcal X\times\mathcal Y$, either $|f(x,y)|=1$, in which case we say $f$ is defined at $(x,y)$, or $f(x,y)=\mathcal Z$, in which case we say $f$ is undefined at $(x,y)$. If $f$ is defined at $(x,y)$, we write $f(x,y)=z$ in place of $f(x,y)=\{z\}$, and if $f$ is undefined at $(x,y)$, we write $f(x,y)=*$, instead of $f(x,y)=\mathcal Z$.
\begin{ntdefinition}[Author Definition23: $P(F)$]
Given a relation $F\subseteq\mathcal X\times\mathcal Y\times\mathcal Z$, a partition of $F$ is a family $\Pi$ of monochromatic rectangles for $F$, such that every $(x,y)\in\mathcal X\times\mathcal Y$ appears in exactly one rectangle of $\Pi$. The partition number of a relation $F\subseteq\mathcal X\times\mathcal Y\times\mathcal Z$, denoted $P(F)$ is the smallest possible size of a partition of $F$.
\end{ntdefinition}
\begin{ntdefinition}[Author Definition24: $C(F)$]
Given a relation $F\subseteq\mathcal X\times\mathcal Y\times\mathcal Z$, a cover of $F$ is a family $\Gamma$ of monochromatic rectangles for $F$, such that every $(x,y)\in\mathcal X\times\mathcal Y$ appears in at least one rectangle of $\Gamma$. The cover number of a relation $F\subseteq\mathcal X\times\mathcal Y\times\mathcal Z$, denoted $C(F)$ is the smallest possible size of a cover of $F$.
\end{ntdefinition}
In the context of Boolean-valued functions, the usual notion of cover number (see~\cite{ref42} Section4.2,p97), and related notion of non-deterministic communication complexity, only require that the 1s of the communication matrix are covered.
\begin{ntdefinition}[Author Definition25: $C_1(F)$]
Given a partial two-player function $f:\mathcal X\times\mathcal Y\to\{0,1\}$, a 1-cover of $F$ is a family $\Gamma$ of 1-monochromatic rectangles for $F$, such that every $(x,y)\in\mathcal X\times\mathcal Y$ having $f(x,y)=1$ appears in at least one rectangle of $\Gamma$. The1-cover number of $f$, denoted $C_1(F)$ is the smallest possible size of a 1-cover of $F$. We then define the non-deterministic communication complexity of $f$ to be $N(f)=\lceil\log C_1(f)\rceil$.
\end{ntdefinition}
\begin{ntdefinition}[Author Definition26]
A protocol $\pi$ over $\mathcal X\times\mathcal Y$ is a rooted tree:
\begin{itemize}
\item Each node $v$ is associated with a rectangle $\pi^{-1}(v)=A\times B\subseteq\mathcal X\times\mathcal Y$.
\item Each non-leaf node $v$, with $\pi^{-1}(v)=A\times B$, is labeled by either(a) a partition $A=A_0\mathbin{\dot\cup}A_1$ of $A$, in which case we say it is Alice's node or(b) a partition $B=B_0\mathbin{\dot\cup}B_1$ of $B$, in which case we say it is Bob's node.
\item The rectangle associated with the root is $\mathcal X\times\mathcal Y$.
\item If a non-leaf node $v$ of Alice has $\pi^{-1}(v)=A\times B$ and is labeled by a partition $A=A_0\mathbin{\dot\cup}A_1$ of $A$, then for each $c\in\{0,1\}$ there will be a unique child $v_c$ of $v$, with $\pi^{-1}(v_c)=A_c\times B$; similarly for Bob's nodes.
\end{itemize}
\end{ntdefinition}
\begin{ntdefinition}[Author Definition27: $L(F),D(F)$]
Let $F\subseteq\mathcal X\times\mathcal Y\times\mathcal Z$ be a relation, and $\pi$ a protocol over $\mathcal X\times\mathcal Y$. We say that $\pi$ is a protocol for $F$ if, for any leaf $\ell$ of $\pi$, the rectangle $\pi^{-1}(\ell)$ is monochromatic for $F$. We then let $L(F)$ be the smallest possible number of leaves in a protocol for $F$, and we let $D(F)$ be the smallest possible depth of a protocol for $F$.
\end{ntdefinition}
Note that the rectangles associated with the leaves of a protocol for $F$ form a partition of $F$, any partition of $F$ is a cover of $F$, and any cover of $F$ contains a 1-cover of $F$, and hence we get the following:
\begin{ntlemma}[Author Lemma28]
For any relation $F\subseteq\mathcal X\times\mathcal Y\times\mathcal Z$, $L(F)\geq P(F)\geq C(F)$, and for a partial two-player function $f:\mathcal X\times\mathcal Y\to\{0,1\}$, $C(f)\geq C_1(f)$.
\end{ntlemma}
\begin{ntdefinition}[Author Definition29]
A DAG-like protocol $\gamma$ over $\mathcal X\times\mathcal Y$ is a directed acyclic graph:
\begin{itemize}
\item $\gamma$ has a single source node, called the root, and one or more sink nodes, called leaves.
\item Each node $v$ is associated with a rectangle $\gamma^{-1}(v)=A\times B\subseteq\mathcal X\times\mathcal Y$.
\item Each non-leaf node $v$, with $\gamma^{-1}(v)=A\times B$, is labeled by either(a) a partition $A=A_0\mathbin{\dot\cup}A_1$ of $A$, in which case we say it is Alice's node or(b) a partition $B=B_0\mathbin{\dot\cup}B_1$ of $B$, in which case we say it is Bob's node.
\item The rectangle associated with the root is $\mathcal X\times\mathcal Y$.
\item If a non-leaf node $v$ of Alice has $\gamma^{-1}(v)=A\times B$ and is labeled by a partition $A=A_0\mathbin{\dot\cup}A_1$ of $A$, then for each $c\in\{0,1\}$ there will be an edge from $v$ to a node $v_c$ of $\gamma$, which must be such that $A_c\times B\subseteq\gamma^{-1}(v_c)$; similarly for Bob's nodes.
\end{itemize}
\end{ntdefinition}
\begin{ntdefinition}[Author Definition30: $S(F)$]
Let $F\subseteq\mathcal X\times\mathcal Y\times\mathcal Z$ be a relation, and $\gamma$ a DAG-like protocol over $\mathcal X\times\mathcal Y$. We say that $\gamma$ is a DAG-like protocol for $F$ if, for any leaf $\ell$ of $\gamma$, the rectangle $\gamma^{-1}(\ell)$ is monochromatic for $F$. We then let $S(F)$ be the smallest possible number of nodes in a DAG-like protocol for $F$.
\end{ntdefinition}
Note that:(1) a protocol with $k$ leaves is a DAG-like protocol with $2k-1$ nodes,(2) the rectangles associated with the leaves of a DAG-like protocol form a cover of $F$, and(3) any DAG with 1 source node,$k$ sink nodes, and maximum out-degree 2 has at least $2k-1$ nodes in total. Hence:
\begin{ntlemma}[Author Lemma31]
$S(F)\leq2L(F)-1$ and $S(F)\geq2C(F)-1$.
\end{ntlemma}
